\chapter{ROS2-Implementierung mit Python}

\begin{lernziele}
Nach diesem Kapitel verstehen Sie den Aufbau eines einfachen ROS2-Python-Nodes.
\end{lernziele}

\section{Package erstellen}
Ein Python-Package wird typischerweise mit folgendem Befehl erstellt:

\begin{lstlisting}[language=bash]
ros2 pkg create robot_basic_nodes --build-type ament_python --dependencies rclpy std_msgs sensor_msgs geometry_msgs
\end{lstlisting}

\section{Ein einfacher Node}
\begin{lstlisting}[language=Python]
import rclpy
from rclpy.node import Node
from std_msgs.msg import String

class StatusNode(Node):
    def __init__(self):
        super().__init__('status_node')
        self.publisher = self.create_publisher(String, '/robot/status', 10)
        self.timer = self.create_timer(1.0, self.publish_status)

    def publish_status(self):
        msg = String()
        msg.data = 'Robot is running'
        self.publisher.publish(msg)

def main():
    rclpy.init()
    node = StatusNode()
    rclpy.spin(node)
    node.destroy_node()
    rclpy.shutdown()
\end{lstlisting}

\section{Subscriber}
Ein Subscriber empfängt Nachrichten von einem Topic. Der ObjectDetector würde zum Beispiel \texttt{/camera/image\_raw} abonnieren.

\section{Build und Start}
\begin{lstlisting}[language=bash]
colcon build
source install/setup.bash
ros2 run robot_basic_nodes status_node
\end{lstlisting}

\begin{hinweisbox}
Der Code ist bewusst minimal. In realen Projekten kommen Fehlerbehandlung, Parameter, Logging und Launch-Dateien hinzu.
\end{hinweisbox}

\begin{uebungbox}
Implementieren Sie einen \texttt{battery\_monitor\_node}, der regelmäßig einen simulierten Akkustand veröffentlicht.
\end{uebungbox}
